\subsection{Permanencia de propiedades conjuntistas y topológicas de morfismos
por descenso fielmente plano.}
\label{subsection:IV.2.6-fr}

\begin{proposition}[2.6.1]
\EGAReaderTarget{ega4.c04.prop.2-6-1}{IV.2.6.1-fr}
Sean $f:X\to Y$ un $S$-morfismo de $S$-preesquemas, $g:S'\to S$ un
morfismo sobreyectivo, $X'\to X_{(S')}$, $Y'=Y_{(S')}$,
$f'=f_{(S')}:X'\to Y'$. Consideremos, para un morfismo, la propiedad de ser:
\begin{enumerate}
\item[\textup{(i)}] sobreyectivo;
\item[\textup{(ii)}] inyectivo;
\item[\textup{(iii)}] de fibras finitas \textup{(como \emph{conjuntos})};
\item[\textup{(iv)}] biyectivo;
\item[\textup{(v)}] radicial.
\end{enumerate}
Entonces, si $\boldsymbol{P}$ designa una de las propiedades precedentes y si $f'$ posee la
propiedad $\boldsymbol{P}$, lo mismo ocurre con $f$.
\end{proposition}

Como el morfismo de proyección $Y'\to Y$ también es sobreyectivo (\textbf I, 3.5.2),
es posible, en virtud de (\textbf I, 3.3.11), limitarse al caso en que $Y=S$,
$Y'=S'$. Para todo $y\in Y$ (resp. $y'\in Y'$), denotemos por $X_y$ (resp.
$X'_{y'}$) el preesquema fibra $f^{-1}(y)$ (resp. $f'^{-1}(y)$) (\textbf I, 3.6.2);
se sabe que, para $y'\in Y'$, $y=g(y')$, se tiene un isomorfismo canónico $X'_{y'}\xrightarrow{\sim}X_y\otimes_{\mathbf{k}(y)}\boldsymbol{k}(y')$ (\textbf{I}, 3.6.4);
como el morfismo $\operatorname{Spec}(\boldsymbol{k}(y'))\to\operatorname{Spec}(\boldsymbol{k}(y))$
es sobreyectivo, lo mismo ocurre con la proyección $X'_{y'}\to X_y$
(\textbf I, 3.5.2); por tanto, si $X'_{y'}$ no es vacío (resp. tiene a lo sumo un
punto, resp. es un conjunto finito), lo mismo ocurre con $X_y$; como
$g:Y'\to Y$ es sobreyectivo, esto demuestra la proposición en los casos (i),
(ii) y (iii), y (iv) resulta de (i) y (ii); por último, para demostrar (v),
basta mostrar que, si $f'$ es universalmente inyectivo (\textbf I, 3.5.11), lo mismo
ocurre con $f$; ahora bien, sea $Y_1\to Y$ un morfismo arbitrario; pongamos
$X_1=X\times_Y Y_1$ y $f_1=f_{(Y_1)}$. Por otra parte, pongamos $Y'_1=Y_1\times_Y Y'$, $X'_1=X'\times_{Y'}Y'_1=X_1\times_{Y_1}Y'_1$, $f'_1=f'_{(Y'_1)}=(f_1)_{(Y'_1)}:X'_1\to Y'_1$;
como $g'=g_{(Y_1)}:Y'_1\to Y'$ es sobreyectivo (\textbf I, 3.5.2) y $f'$ es
universalmente inyectivo, $f'_1$ es inyectivo, y de (ii) resulta que
$f_1$ es inyectivo, de donde se sigue nuestra afirmación.

